% Source: https://artofproblemsolving.com/community/c3h2560522_expected_value_problem_up_vs_down
% Stable identity: aops-topic-2560522-q01
Let $E_k$ be the expected number of steps from $k$. These are finite: from any position, the next $N$ steps all going down reach $0$, so each block of $N$ steps ends the process with probability at least $2^{-N}$, and the number of steps has a geometric tail.

Clearly $E_0=E_N=0$, and conditioning on the first step,
\[E_k=1+\dfrac12E_{k-1}+\dfrac12E_{k+1}\qquad(0<k<N).\]
The function $f(k)=k(N-k)$ satisfies the same equation: $f(k-1)+f(k+1)=2k(N-k)-2$, so $1+\dfrac12\left(f(k-1)+f(k+1)\right)=f(k)$. Hence $D_k=E_k-f(k)$ satisfies $D_k=\dfrac12\left(D_{k-1}+D_{k+1}\right)$, i.e. $D_{k+1}-D_k=D_k-D_{k-1}$: the $D_k$ form an arithmetic sequence. Since $D_0=D_N=0$, all $D_k=0$, so $E_k=k(N-k)$. $\blacksquare$

\textbf{(a)} With $k=500$ and $N=1000$, the expected number of steps is $500\cdot500=\boxed{250000}$.
